\textbf{Задача 5.}
Докажите, что для любой последовательности $\{a_i\}, a_i \in \{0, 1\}$ существует последовательность строго вложенных языков $A_0 \subsetneq A_1 \subsetneq \dots$, такая что $A_i \in \mathbf{P}$ при $a_i = 0$ и $A_i$ является $\mathbf{NP}$-полным при $a_i = 1$.

\begin{proof}
Заметим, что задача поиска гамильтонова цикла является $\mathbf{NP}$-полной задачей, а нахождение по $N$ целой части $N^{\frac{1}{c}}$ для некоторой константы $c$ лежит в $\mathbf{P}$ (бинарный поиск). Также известно, что $S(N, K) = N^{\frac{1}{2}} + \dots + N^{\frac{1}{K}} < N$ при достаточно больших $N$ и $\forall$ фиксированном $K \in \N$. Определим язык $A_i, i \ge 1$ так: это язык, состоящий из графов, таких, что на первых $S(N, i - 1)$ вершинах граф произволен, последние $N - S(N, i)$ вершин не лежат ни в каком ребре, а на остальных $n^{\frac{1}{i}}$ вершинах граф произволен при $a_i = 0$ или имеет гамильтонов цикл (либо не имеет ребер) при $a_i = 1$. Тут $N$ - число вершин графа. Также будем предполагать, что если $N \leq S(N, i)$, то такой граф лежит в языке.
\end{proof}

\textbf{Альтернативная идея:} Сначала построим последовательность чередующихся языков \p и \npc, а затем при необходимости выкинем ненужные. 

Рассмотрим $A_i$ = \{не делящиеся на $2^i$ натуральные числа\}. Возьмём $A_i \subset B_i \subset A_{i+1}$ так, чтобы всё выполнялось

Из $A_{i+1} \setminus A_i$ возьмём что-нибудь \np-полное. Можно, например, взять любой \np-полный язык, домножить на $2^{i+1}$ и прибавить $2^i$. Получится хрень, не делящаяся на $2^{i+1}$, но делящаяся на $2^{i}$, то есть из $A_{i+1} \setminus A_i$.

Тогда к этому делу добавим $A_i$ и получится \np-полный язык, потому что его распознавание полиномиально эквивалентно распознаванию исходного \np-полного языка.